\documentclass[11pt,a4paper]{article}

% ---- Packages ----
\usepackage[utf8]{inputenc}
\usepackage[T1]{fontenc}
\usepackage{amsmath,amssymb,amsthm}
\usepackage{mathtools}
\usepackage{geometry}
\geometry{margin=2.4cm}
\usepackage{enumitem}
\usepackage{booktabs}
\usepackage{hyperref}
\hypersetup{colorlinks=true,linkcolor=blue!60!black,citecolor=blue!60!black,urlcolor=blue!60!black}
\usepackage{xcolor}
% \usepackage{natbib}
% \bibliographystyle{plainnat}

% ---- Theorem environments ----
\theoremstyle{proposition}
\newtheorem{theorem}{Theorem}[section]
\newtheorem{proposition}[theorem]{Proposition}
\newtheorem{lemma}[theorem]{Lemma}
\newtheorem{corollary}[theorem]{Corollary}
\newtheorem{remark}[theorem]{Remark}
\newtheorem{definition}[theorem]{Definition}
\newtheorem{assumption}[theorem]{Assumption}
\newtheorem{example}[theorem]{Example}
% \theoremstyle{definition}

% ---- Shortcuts ----
\newcommand{\R}{\mathbb{R}}
\newcommand{\CoP}{\mathrm{CoP}}
\newcommand{\tr}{\operatorname{tr}}
\newcommand{\diag}{\operatorname{diag}}
\newcommand{\LS}{\mathcal{L}}
\newcommand{\comment}{\textcolor{red}}

% ---- Title ----
\title{\textbf{Torsional friction in planar contact}}
\author{Sébastien Kleff\\[2pt]
\small\texttt{sk8001@nyu.edu}}
\date{\today}

\begin{document}
\maketitle


\begin{abstract}
Tactile sensors are widely used to detect slip, but much less is known about what can actually be inferred about frictional contact limits from tactile measurements. This question matters because detecting that slip is occurring or imminent is not the same as knowing which tangential and torsional loads the contact can still support. That distinction is important for control design. In this report, we study this question for a single planar contact patch under quasi-static Coulomb friction. We first show that coarse normal force information, consisting of the total normal force and center of pressure, is generally insufficient to determine the admissible tangential wrench set: two contacts can have the same total normal force and CoP while supporting different combinations of tangential force and torsional moment. We then study what can still be inferred from partial measurements. In particular, we derive upper and lower bounds on the admissible wrench set from coarse measurements and from spatially distributed taxel-wise normal forces. The latter provide tighter bounds by reducing the ambiguity on how the normal force is distributed over the contact patch.
\end{abstract}

\section{Introduction}
Modern tactile sensors provide rich local information about contact pressure, shape, and deformation. Yet tactile features used for control are often selected mainly from empirical performance, for example contact pose, resultant force, image features, or local pressure measurements. A more basic question is less often made explicit: given a particular tactile measurement, what can actually be inferred about the physical interaction at the contact?
We study this question in the setting of a rigid planar contact patch under Coulomb friction. Rather than trying to characterize every possible tactile signal or manipulation task, we focus on one concrete control-relevant quantity: the set of tangential forces and torsional moments that the contact can support without violating the friction model.
The key difficulty is that this admissible wrench set depends on how the normal load is distributed over the patch. Coarse measurements such as the total normal force and center of pressure constrain this load distribution, but do not determine it uniquely. Spatial tactile measurements provide additional information and therefore reduce this ambiguity.
The objective of this report is to make this relationship explicit. We ask \textbf{how the admissible frictional wrench set can be bounded from different levels of normal force information}, from total force and CoP to spatially distributed taxel-wise forces, \textbf{without requiring the full normal pressure field}.

% \section{Result}
\section{Preliminaries (measures)}
\label{sec:measure-preliminaries}

We briefly introduce the measure-theoretic notation used in the following.
Only the notions needed for the contact model are recalled here; standard reference
is \cite{evans2015measure}.

\subsection{Force as a measure.}
Let $P\subset\mathbb R^2$ denote the contact patch, equipped with its Borel
$\sigma$-algebra $\mathcal B(P)$. A \textit{nonnegative measure} $\nu_N:\mathcal B(P)\rightarrow[0,+\infty]$ assigns to every measurable subset $A\subseteq P$ a nonnegative scalar $\nu_N(A)$, and is countably additive, i.e. for any pairwise disjoint measurable sets $A_i \subseteq P$,
\begin{align}
\nu_N\left(\bigcup_i A_i\right) =\sum_i \nu_N(A_i).
\end{align}
In the present context, $\nu_N$ represents the (spatially distributed) normal force, i.e. $\nu_N(A)$ is the normal force carried by the region A. Hence $\nu_N(A)$ has units of force (N), and the resultant normal force is $f_N=\nu_N(P)$. We restrict attention to \textit{finite nonnegative Radon measures}, denoted $\nu_N\in\mathcal M_+(P)$. Informally, a Radon measure is a "well-behaved" measure. In particular, on the compact sets considered here, it is finite and its value on measurable sets can be approximated from compact subsets. For our purpose, it essentially means that our measure does not assign an infinite force to $P$ and it can represent either the force induced by an ordinary pressure field or an idealized concentrated point/line load.

We represent the tangential force by an $\mathbb R^2$-valued finite Radon measure $\nu_T=(\nu_{T,1},\nu_{T,2})$ where each component $\nu_{T,i}$ is a finite signed Radon measure on $P$. Thus, for every measurable region $E\subseteq P$, the resultant tangential force carried by $E$ is
\begin{align}
    \nu_T(E)
    =
    \begin{bmatrix}
    \nu_{T,1}(E)\\
    \nu_{T,2}(E)
    \end{bmatrix}
    \in\mathbb R^2,
\end{align}

\subsection{Area measure and pressure fields.}
Let $dA$ denote the usual area element on the planar patch $P$.
A \textit{normal pressure field} is a nonnegative function $\rho_N:P\rightarrow\mathbb R_+$ with units of force per unit area. If $\rho_N\in L^1_+(P)$, which simply means that $\rho_N$ has finite integral over $P$
\begin{align}
\rho_N(x)\ge 0\quad\text{a.e.},
\qquad
\int_P \rho_N(x)\,dA<\infty.
\end{align}
then it defines a finite normal force measure: the force carried by any region $E\subseteq P$ is given by
\begin{align}
\nu_N(E) = \int_E \rho_N(x)\,dA.
\end{align}
which we also write equivalently as $d\nu_N(x)=\rho_N(x)\,dA$. In that case, $\rho_N$ is called the \textit{density} of the normal force measure $\nu_N$ with respect to area\footnote{More generally, the Radon-Nikodym theorem guarantees the existence of a density for measures under suitable assumptions, such as "absolute continuity" which means physically that a subset of zero area (e.g. a point or a line) carries zero normal force. See~\cite{evans2015measure}, Theorem 1.30, p. 50}. Not every normal force measure admits such a pressure field. For example the Dirac measure $\nu_N=f_N\delta_c$ represents an idealized normal force $f_N$ concentrated at the single point $c$. It cannot be represented by an ordinary pressure field over area. In this note we adopt the measure-theoretic language (instead of only $L_+^1$ densities) so that we cover both "classical" pressure fields and idealized concentrated loads.

This notation allows ordinary pressure fields and concentrated loads to be handled by the same integration equations. For a measurable scalar or vector function $g$, the integral 
\begin{align}
\int_P g(x)\,d\nu_N(x)
\end{align}
can be interpreted as a force-weighted spatial integral. 
Whenever $\nu_N$ admits a normal pressure field $\rho_N$,
\begin{align}
\int_P g(x)\,d\nu_N(x) = \int_P g(x)\rho_N(x)\,dA.
\end{align} 
For a Dirac measure,
\begin{align}
\int_P g(x)\,d(f_N\delta_c)(x) = f_Ng(c).
\end{align}

\subsection{Moments of the normal force measure.}
The zeroth moment of $\nu_N$ is its total measure,
\begin{align}
m_0(\nu_N) \triangleq \int_P 1\,d\nu_N = \nu_N(P) = f_N.
\end{align}
In the present mechanical interpretation, this is the \textit{resultant normal force}. The first moment is
\begin{align}
m_1(\nu_N) \triangleq \int_P x\,d\nu_N(x)\in\mathbb R^2.
\end{align}
If $f_N>0$, the \textit{center of pressure} is the normalized first moment, 
\begin{align} 
c[\nu_N] &\triangleq \frac{m_1(\nu_N)}{m_0(\nu_N)} \\
         & = \frac{1}{f_N} \int_P x\,d\nu_N(x).
\end{align}
Hence knowing the total normal force and the center of pressure is equivalent to knowing the zeroth and first moments of the normal force measure $m_0(\nu_N)=f_N, m_1(\nu_N)=f_Nc[\nu_N]$. 

\section{Problem formulation}
\subsection{Planar contact model}

Assume that the planar contact patch $P\subset\mathbb R^2$ is compact and convex,
and let $\nu_N\in\mathcal M_+(P)$ and $\nu_T$ be the normal and tangential force
measures introduced above. We assume
\begin{align}
f_N=m_0(\nu_N)>0,
\end{align}
so that the center of pressure $c$ is well-defined. Since $P$ is convex and
$\nu_N$ is nonnegative, $c\in P$.
We further assume a known and spatially uniform Coulomb friction coefficient
$\mu>0$.
We impose Coulomb friction in measure form by requiring that there exists a
measurable density $\gamma_T:P\rightarrow\mathbb R^2$ of $\nu_T$ with respect
to $\nu_N$ such that
\begin{align}
\label{eq:coulomb_measure}
d\nu_T(x)=\gamma_T(x)\,d\nu_N(x),
\qquad
\|\gamma_T(x)\|_2\le\mu
\quad \nu_N\text{-a.e. on }P.
\end{align}
We call $\gamma_T$ the \textit{friction ratio field}. It can be interpreted as
the tangential force per unit normal force. 
The above condition expresses both that tangential friction force can only be carried where normal force is present
and that its local magnitude is bounded by $\mu$\footnote{In particular, $d\nu_T=\gamma_T\,d\nu_N$ implies
$\nu_T\ll\nu_N$.}. When $\nu_N$ admits the normal pressure field $\rho_N$ introduced above, i.e. $d\nu_N=\rho_N\,dA$, the tangential force measure also admits a density with respect to area,
\begin{align}
\rho_T(x)=\gamma_T(x)\rho_N(x),
\end{align}
called the \textit{tangential traction field}. The measure-form Coulomb condition then reduces to the classical local condition
\begin{align}
\label{eq:coulomb_local}
\|\rho_T(x)\|_2\le\mu\rho_N(x)
\qquad \text{a.e. on }P.
\end{align}
We define the \textit{resultant tangential force} and \textit{resultant torsional moment} about the center of pressure $c$ as
\begin{align}
f_T \triangleq \nu_T(P)
= \int_P\gamma_T(x)\,d\nu_N(x),
\qquad
\tau_N \triangleq
\int_P (x-c)\times\gamma_T(x)\,d\nu_N(x).
\end{align}
We finally define the skew-symmetric matrix
\begin{align}
J \triangleq
\begin{bmatrix}
0&-1\\
1&0
\end{bmatrix},
\end{align}
so that for $x,y\in\mathbb R^2$ $x\times y=(Jx)^\top y$

\begin{definition}[Admissible tangential wrench set]
    Define a normal force measure $\nu_N\in\mathcal M_+(P)$. We define \textit{admissible tangential wrench set} :
    \begin{align}
    \mathcal W(\nu_N)
    \triangleq
    \left\{
    (f_T,\tau_N)\in\mathbb R^2\times\mathbb R:
    \begin{array}{l}
    \exists\,\gamma_T:P\rightarrow\mathbb R^2
    \text{ measurable such that}\\
    \|\gamma_T(x)\|_2\le\mu
    \quad\nu_N\text{-a.e.},\\
    f_T=\displaystyle\int_P\gamma_T(x)\,d\nu_N(x),\\
    \tau_N=\displaystyle\int_P
    (x-c[\nu_N])\times\gamma_T(x)\,d\nu_N(x)
    \end{array}
    \right\}.
    \end{align}
    The set $\mathcal W(\nu_N)$ is a non-empty compact subset of $\R^3$ containing all resultant tangential wrenches (tangential force and torsional moment) that are compatible with the normal force measure $\nu_N$ and the local Coulomb friction condition~\eqref{eq:coulomb_measure}.
\end{definition}




\begin{proposition}[Support function of the admissible tangential wrench set]
    \label{prop:support}
    For any $\nu_N\in\mathcal M_+(P)$ and for every $q=(a,b)\in\mathbb R^2\times\mathbb R$, the support function of $\mathcal W(\nu_N)$ is
    \begin{align}
    h_{\mathcal W(\nu_N)}(q) = \mu \int_P \left\|a+bJ(x-c[\nu_N])\right\|_2\,d\nu_N(x).
    \end{align}
\end{proposition}
\begin{proof}
    By definition, $\forall q=(a,b)\in\R^2\times \R$
    \begin{align}
    h_{\mathcal W(\nu_N)}(q) =\sup_{(f_T,\tau_N)\in\mathcal W(\nu_N)}\left(a^\top f_T+b\tau_N\right).
    \end{align}
    For any friction ratio field $\gamma_T$
    \begin{align}
    a^\top f_T+b\tau_N=\int_P\left[a+bJ(x-c[\nu_N])\right]^\top\gamma_T(x)\,d\nu_N(x).
    \end{align}
    Using Cauchy--Schwarz and
    $\|\gamma_T(x)\|_2\le\mu$,
    \begin{align}
    a^\top f_T+b\tau_N\le\mu\int_P\left\|a+bJ(x-c[\nu_N])\right\|_2\,d\nu_N(x).
    \end{align}
    The bound is attained by choosing $\gamma_T(x)$ with magnitude $\mu$ and aligned with
    $a+bJ(x-c[\nu_N])$ whenever this vector is nonzero:
    \begin{align}
        \gamma_T^{\star}(x)= \begin{cases}
                \mu\frac{a + bJ(x-c[\nu_N])}{\|a + bJ(x-c[\nu_N])\|_2}~~&\text{if}~\|a + bJ(x-c[\nu_N])\|_2\neq 0\\
                0~~&\text{otherwise}
        \end{cases} 
    \end{align}
\end{proof}


\subsection{Measurement-conditioned support bounds}
\begin{table}[t]
    \centering
    \small
    \resizebox{\textwidth}{!}{\begin{tabular}{llll}
        \toprule
        Symbol & Meaning & Remark & Lives in \\
        \midrule
        $\nu_N$ 
        & normal force measure
        & true $\nu_N^\star$ is unknown
        & $\mathcal M_+(P)$
        \\
        $c[\nu_N]$
        & CoP induced by measure $\nu_N$
        & depends on $\nu_N$
        & $P\subset\mathbb R^2$
        \\
        $q=(a,b)$
        & wrench support direction
        & prescribed
        & $\mathbb R^2\times\mathbb R$
        \\
        $\mathcal M$
        & measurement operator
        & fixed by the sensing model
        & $\mathcal M_+(P)\rightarrow\mathbb R^r$
        \\
        $y$
        & measurement value 
        & observed data from $(\mathcal M,\nu_N)$
        & $\mathbb R^r$
        \\
        $\Pi_{\mathcal M}(y)$
        & all measures compatible with $y$
        & depends on $(\mathcal M,y)$
        & $\subseteq\mathcal M_+(P)$
        \\
        $\mathcal C_{\mathcal M}(y)$
        & CoPs compatible with $y$
        & depends on $(\mathcal M,y)$
        & $\subseteq P$
        \\
        % $h_{\mathcal W(\nu_N)}(q)$
        % & support of the wrench set for $\nu_N$
        % & depends on $(\nu_N,q)$
        % & $\mathbb R$
        % \\
        $\underline h_{\mathcal M}(q;y)$
        & measurement-conditioned lower bound
        & depends on $(\mathcal M,y,q)$
        & $\mathbb R$
        \\
        $\overline h_{\mathcal M}(q;y)$
        & measurement-conditioned upper bound
        & depends on $(\mathcal M,y,q)$
        & $\mathbb R$
        \\
        \bottomrule
    \end{tabular}}
    \caption{Notations.}
    \label{tab:measurement_notation}
\end{table}

We focus here on the case where we do not know the \textit{true} force measure $\nu_N^\star$ but we have only access to some measurements of it (e.g. resultant normal force, CoP, etc.). 
The goal is to investigate what can be known about the support of $\mathcal W(\nu_N^\star)$ given partial measurements of $\nu_N^\star$. 
\begin{definition}[Measurement operator]\label{def:measurement_operator}
    We define a \textit{measurement operator} as a function $\mathcal M:\mathcal M_+(P)\rightarrow\mathbb R^r$. We assume that its components are linear integral measurements of the form
    \begin{align}
    \mathcal M(\nu_N) = 
    \begin{bmatrix}
    \displaystyle\int_P \phi_1(x)\,d\nu_N(x)\\
    \vdots\\
    \displaystyle\int_P \phi_r(x)\,d\nu_N(x)
    \end{bmatrix},\qquad \forall \nu_N \in \mathcal M_+(P)
    \end{align}
    where the prescribed measurable functions
    $\phi_j:P\rightarrow\mathbb R$ are called the \textit{measurement functions}.
    In the following, we denote a measurement value (or observation) induced by the (unknown) true force measure $\nu_N^\star$ as $y = \mathcal M(\nu_N^\star) \in\mathbb R^r$.
\end{definition}

\begin{definition}[Measurement-compatible normal force family]
    Given a measurement operator $\mathcal M:\mathcal M_+(P)\rightarrow\mathbb R^r$ and the observation $y\in\R^r$, we define the set of \textit{normal force measures that are compatible with observation $y$} 
    \begin{align}
    \Pi_{\mathcal M}(y) \triangleq\left\{\nu_N\in\mathcal M_+(P):\mathcal M(\nu_N)=y\right\}.
    \end{align}
    The set $\Pi_{\mathcal M}(y)$ contains all the normal force measures that are indistinguishable under measurement $y$. Note that this is the pre-image of $y$ through $\mathcal{M}$ : $\Pi_{\mathcal{M}(y)} = \mathcal{M}^{-1}[\{y\}]$.
\end{definition}

\begin{assumption}
    From now, on we assume \(\forall \nu_N\in\Pi_{\mathcal M}(y)\), $\nu_N(P) >0$. Purpose: we restrict the measurement-compatible family to measures that have a well-defined CoP.
    % From now on, we assume that the measurement operator includes at least the zeroth moment, i.e. $r\geq 1$ and $\phi_1(x)=1$. \textcolor{red}{Not true for $\mathcal{M}_F$ !}
\end{assumption}

\begin{definition}[Set of measurement-compatible CoPs]\label{def:CM}
    We define
    \begin{align}
    \mathcal C_{\mathcal M}(y)
    \triangleq
    \left\{
    c[\nu_N]:
    \nu_N\in\Pi_{\mathcal M}(y)
    \right\}
    \end{align}
    as the set of CoPs compatible with the observation $y \in \mathbb{R}^r$.    
\end{definition}

Given an observation $y$ of the unknown true force measure $\nu_N^\star$, we can compute measurement-conditioned lower/upper bounds on the support function $h_{\mathcal W(\nu_N^\star)}$ by optimizing over $\Pi_{\mathcal M}(y)$. 

\begin{definition}[Measurement-conditioned support bounds]\label{def:measurement_support_bounds}
Let $\mathcal M:\mathcal M_+(P)\rightarrow\mathbb R^r$ be a measurement operator as defined above and let $y=\mathcal M(\nu_N^\star)$ denote the observation induced by the unknown true normal force measure $\nu_N^\star$. For a wrench direction $q\in\mathbb R^2\times\mathbb R$, we define the \textit{measurement-conditioned bounds in direction $q$}:
\begin{align}
    \label{eq:def_mea_bounds}
    \underline h_{\mathcal M}(q;y)
    &\triangleq
    \inf_{\nu_N\in\Pi_{\mathcal M}(y)}
    h_{\mathcal W(\nu_N)}(q),\\
    \overline h_{\mathcal M}(q;y)
    &\triangleq
    \sup_{\nu_N\in\Pi_{\mathcal M}(y)}
    h_{\mathcal W(\nu_N)}(q).
\end{align}
\end{definition}

% By this definition, for every $q\in\mathbb R^2\times\mathbb R$ and every
% $\nu_N\in\Pi_{\mathcal M}(y)$,
% \begin{align}
%     \underline h_{\mathcal M}(q;y)
%     \le
%     h_{\mathcal W(\nu_N)}(q)
%     \le
%     \overline h_{\mathcal M}(q;y).
% \end{align}

\begin{remark}[How to use the bounds $\underline h_{\mathcal M}(q;y), \overline h_{\mathcal M}(q;y)$]
    A practical consequence of knowing the bounds is that we can test the slip feasibility of a candidate tangential wrench $w$: guaranteed feasible if $\forall q,~\langle q,w\rangle \le \underline h_{\mathcal M}(q;y)$.
    Conversely, if $\exists q,~\langle q,w\rangle >\overline h_{\mathcal M}(q;y)$, then $w$ is infeasible for every normal force measure compatible with the measurement $y$\footnote{Because all support functions are positively homogeneous in $q$, it is sufficient in practice to search over one normalized representative of each nonzero ray. After fixing a consistent scaling between force and torque coordinates, one may for instance impose $\|q\|_2=1$.}.
\end{remark}


\begin{remark}[Relation with inner/outer wrench sets]
    We can relate the bounds of Definition~\ref{def:measurement_support_bounds} to outer/inner approximations of the admissible wrench set
    \begin{align}
        \mathcal W^+(y)
        &\triangleq
        \bigcup_{\nu_N\in\Pi_{\mathcal M}(y)}\mathcal W(\nu_N),
        &
        \mathcal W^-(y)
        &\triangleq
        \bigcap_{\nu_N\in\Pi_{\mathcal M}(y)}\mathcal W(\nu_N).
        \end{align}
    Then $h_{\mathcal W^+(y)}(q) =\overline h_{\mathcal M}(q;y)$ and $h_{\mathcal W^-(y)}(q) \le\underline h_{\mathcal M}(q;y)$ (equality need not hold for the inner set). Nevertheless, the lower bound characterizes the inner set exactly through
    \begin{align}
        \mathcal W^-(y)
        =
        \left\{
        w:\,
        \langle q,w\rangle
        \le
        \underline h_{\mathcal M}(q;y)
        \quad \forall q
        \right\}.
    \end{align}
\end{remark}
% \begin{proof}
%     TODO
% \end{proof}

Computing the bounds of Definition~\ref{def:measurement_support_bounds} is difficult in general as it leads to a nonlinear optimization problem over measurement-compatible measures:
\begin{align}
    \underline h_{\mathcal M}(q;y)
    &=
    \inf_{\substack{\nu_N\in\Pi_{\mathcal M}(y)}} %\\ \nu_N(P)>0}}
    \mu
    \int_P
    \left\|
    a+bJ\left(
    x-
    c[\nu_N]
    \right)
    \right\|_2
    \,d\nu_N(x),
    \label{eq:general_nonlinear_lower}
    \\
    \overline h_{\mathcal M}(q;y)
    &=
    \sup_{\substack{\nu_N\in\Pi_{\mathcal M}(y)}} %\\ \nu_N(P)>0}}
    \mu
    \int_P
    \left\|
    a+bJ\left(
    x-
    c[\nu_N]
    \right)
    \right\|_2
    \,d\nu_N(x).
    \label{eq:general_nonlinear_upper}
\end{align}
The objective is nonlinear in the measure $\nu_N$ since the integrand depends
on it through $c[\nu_N]$. Hence, this is not a standard Generalized
Moment Problem (GMP)~\cite{lasserre2010moments}. 
% \eqref{eq:general_nonlinear_lower}--\eqref{eq:general_nonlinear_upper} are equivalent to bi-level problems with outer finite-dimensional optimization over the CoP and an inner GMP:
% \begin{align}
%     \underline h_{\mathcal M}(q;y)
%     =
%     \inf_{c\in\mathcal C_{\mathcal M}(y)}
%     \;
%     \inf_{\substack{
%     \nu_N\in\Pi_{\mathcal M}(y)\\
%     % \nu_N(P)>0\\
%     \int_P(x-c)\,d\nu_N=0}}
%     \mu\int_P
%     \|a+bJ(x-c)\|_2\,d\nu_N(x),
%     \label{eq:parametric_gmp_lower}
% \end{align}
% \begin{align}
%     \overline h_{\mathcal M}(q;y)
%     =
%     \sup_{c\in\mathcal C_{\mathcal M}(y)}
%     \;
%     \sup_{\substack{
%     \nu_N\in\Pi_{\mathcal M}(y)\\
%     % \nu_N(P)>0\\
%     \int_P(x-c)\,d\nu_N=0}}
%     \mu\int_P
%     \|a+bJ(x-c)\|_2\,d\nu_N(x),
%     \label{eq:parametric_gmp_upper}
% \end{align}
% For a fixed $c$, the integrand is fixed and the additional CoP consistency constraint $\int_P(x-c)\,d\nu_N=0$ is linear in $\nu_N$. The inner problem is therefore a GMP. 
The information contained in $\mathcal M$ affects the structure of $\Pi_{\mathcal M}(y)$, which affects the problem's structure. We next specialize problems \eqref{eq:general_nonlinear_lower}, \eqref{eq:general_nonlinear_upper} to particular choices of measurement operators $\mathcal{M}$ in Sections~\ref{sec:M0},~\ref{sec:MF} and \ref{sec:MT}.

% \begin{remark}[Special cases]
%     If the observed data $y$ determine a unique CoP $c$, i.e.
%     $\mathcal C_{\mathcal M}(y)=\{c\}$, then the outer optimization disappears.
%     The bounds \eqref{eq:parametric_gmp_lower},~\eqref{eq:parametric_gmp_upper}
%     reduce to the GMPs
%     \begin{align}
%         \label{eq:gmp_mea_bounds}
%         \underline h_{\mathcal M}(q;y)
%         &=
%         \inf_{\substack{
%         \nu_N\in\Pi_{\mathcal M}(y)}}
%         % \int_P(x-c)\,d\nu_N=0}}
%         \mu\int_P
%         \left\|a+bJ(x-c)\right\|_2\,d\nu_N(x),\\
%         \overline h_{\mathcal M}(q;y)
%         &=
%         \sup_{\substack{
%         \nu_N\in\Pi_{\mathcal M}(y)}}
%         % \int_P(x-c)\,d\nu_N=0}}
%         \mu\int_P
%         \left\|a+bJ(x-c)\right\|_2\,d\nu_N(x).
%     \end{align}
%     The CoP is therefore no longer an optimization variable, but the CoP consistency constraint $\int_P(x-c)\,d\nu_N=0$ remains (inside $\Pi_{\mathcal M}(y)$ since $\mathcal{M}$ imposes $c[\nu_N] = c$ by assumption).
%     \begin{itemize}
%         \item If $\mathcal M$ contains only the resultant normal force and CoP constraints, the above GMP reduces to the coarse-measurement problem studied below \footnote{The corresponding upper GMP coincides, up to normalization and wrench conventions, with the support-function problem used
%         in~[Yann] to construct the De Saxc\'e term for frictional patch
%         contact: both optimize over all nonnegative normal-force
%         distributions having prescribed total force and CoP. Our lower
%         bound addresses the different question of the frictional capacity
%         guaranteed for every normal force distribution compatible with
%         these same measurements.}.

%         \item If $\mathcal M$ additionally contains taxel force constraints $\nu_N(P_i)=F_i$, the CoP constraint couples the otherwise disjoint taxel constraints. This gives the combined coarse-and-taxel
%         measurement case studied below; in particular, we are interested in how the resulting bounds depend on the number and spatial resolution of the taxels.
%     \end{itemize}

%     If the observed data $y$ do not uniquely determine the CoP, then $\mathcal C_{\mathcal M}(y)$ is not a singleton and \eqref{eq:parametric_gmp_lower}--\eqref{eq:parametric_gmp_upper} retain a bi-level structure. In particular, if $\mathcal M$ contains only the taxel-force constraints $\nu_N(P_i)=F_i$, then $\mathcal C_{\mathcal M}(y)$ represents all CoPs compatible with the measured taxel resultants. Characterizing this set, its evolution under taxel refinement, and the resulting outer optimization becomes part of the problem.
% \end{remark}



\begin{table}[t]
    \centering
    \small
    \resizebox{\textwidth}{!}{\begin{tabular}{llll}
        \toprule
        Measurement & Observed data & Compatible CoP set & Resulting problem \\
        \midrule
        $\mathcal M_0$
        &
        $y_0=(f_N,f_Nc)$
        &
        $\mathcal C_{\mathcal M_0}(y_0)=\{c\}$
        &
        GMP (Section~\ref{sec:M0})
        \\[1mm]

        $\mathcal M_F$
        &
        $y_F=(F_1,\ldots,F_r)$
        &
        $\displaystyle
        \mathcal C_{\mathcal M_F}(y_F)
        =
        \left\{
        c[\nu_N]:
        \nu_N(P_i)=F_i,\ \forall i
        \right\}$
        &
        Bi-level (Section~\ref{sec:MF})
        \\[1mm] 
        
        $\mathcal M_T$
        &
        $\displaystyle
        y_T=
        \bigl(y_0,y_F\bigr)$
        &
        $\mathcal C_{\mathcal M_T}(y_T)=\{c\}$
        &
        GMP (Section~\ref{sec:MT})
        \\
        \bottomrule
    \end{tabular}}
    \caption{Measurement operators and corresponding compatible CoP sets.}
    \label{tab:measurement_cases}
\end{table}

\section{$\mathcal{M}_0$ case (coarse measurement)}
\label{sec:M0}

We focus here on a specific case where the available measurements only consist of the resultant normal force $f_N$ and CoP $c$. We show that this partial information is not enough to fully determine the admissible tangential wrench set defined earlier. We do it by providing a counter example in which $f_N,c$ are identical but the torsional friction limits are different.

\begin{definition}[Coarse normal force measurement]\label{def:coarse_measurement}
    Given the true (unknown) normal force measure $\nu_N^\star\in\mathcal M_+(P)$,
    % , let $f_N\triangleq\int_P d\nu_N(x)$ and assume $f_N>0$. 
    we define the coarse measurement as its zeroth and first moments 
    \begin{align}
    y_0 = \mathcal M_0(\nu_N^\star) \triangleq \left(\int_P d\nu_N^\star(x),\;\int_P x\,d\nu_N^\star(x)\right) = (f_N,f_Nc[\nu_N^\star])
    \end{align}
    % where $c\triangleq\frac{1}{f_N}\int_Px\,d\nu_N(x)$ is the CoP.
\end{definition}

\begin{remark}
    The coarse measurement $y_0=\mathcal M_0(\nu_N^\star)$ is a particular case of Definition~\ref{def:measurement_operator}, that extracts the zeroth and first moments through measurement functions:
    \begin{align}
        \phi_1(x)=1,\qquad
        \phi_2(x)=x_1,\qquad
        \phi_3(x)=x_2.
    \end{align}
    % The set of compatible CoPs is a singleton $\mathcal C_{\mathcal M_0}(y_0)=\{c[\nu_N^\star]\}$.
\end{remark}

% \begin{remark} 
%     The GMPs \eqref{eq:gmp_mea_bounds_lower}, \eqref{eq:gmp_mea_bounds_upper} coincide, up to normalization and wrench conventions, with the support-function problem used in~[Yann] to construct the De Saxc\'e term for frictional patch contact: both optimize over all nonnegative normal force distributions having prescribed total force and CoP. Our lower bound addresses the different question of the frictional capacity guaranteed for every normal force distribution compatible with the same measurements.
% \end{remark}        

We now show how to compute the bounds of Definition~\ref{def:measurement_support_bounds} in the special case where only the coarse observation $y_0=\mathcal M_0(\nu_N^\star)$ is available.

\begin{theorem}[Exact coarse lower bound]
    Let $c=c[\nu_N^\star]$. Given the coarse observation $y_0=(f_N, f_Nc)$, for any $q=(a,b)\in\mathbb R^2\times\mathbb R$,
    \begin{align}
            \boxed{\underline h_{\mathcal M_0}(q;y_0) = \mu f_N\|a\|_2}
    \end{align}
\end{theorem}
\begin{proof}
    For every $\nu_N\in\Pi_{\mathcal M_0}(y_0)$,
    Jensen's inequality gives
    \begin{align}
        \int_P\|a+bJ(x-c)\|_2\,d\nu_N(x)
        &\geq
        f_N\left\|
        a+\frac{b}{f_N}J
        \int_P(x-c)\,d\nu_N(x)
        \right\|_2\\
        &=f_N\|a\|_2.
    \end{align}
    Since $P$ is convex, $c\in P$, and the lower bound
    is attained by $\nu_N=f_N\delta_c$.
    % The function $x\mapsto \|a+bJ(x-c)\|_2$ is convex. For any $\nu_N\in\Pi_{\mathcal M_0}(y_0)$, normalizing $\nu_N$ by $f_N$ and applying Jensen's inequality gives
    % \begin{align}
    % \frac{1}{f_N}
    % \int_P
    % \|a+bJ(x-c)\|_2\,d\nu_N(x)
    % &\ge
    % \left\|
    % a+bJ\left(
    % \frac{1}{f_N}\int_P x\,d\nu_N(x)-c
    % \right)
    % \right\|_2\\
    % &=
    % \|a\|_2,
    % \end{align}
    % where the last equality follows from $\mathcal M_0(\nu_N)=y_0$ (which implies $\frac{1}{f_N}\int_P x\,d\nu_N(x)=c$). Therefore,
    % \begin{align}
    % \mu\int_P
    % \|a+bJ(x-c)\|_2\,d\nu_N(x)
    % \ge
    % \mu f_N\|a\|_2
    % \end{align}
    % for every compatible normal force measure. Equality is attained by $\nu_N=f_N\delta_c$, which belongs to $\Pi_{\mathcal M_0}(y_0)$. Hence the infimum is $\mu f_N\|a\|_2$.
\end{proof}

\begin{remark}
As an immediate consequence, setting $a=0$ shows that the coarse measurement does not allow any torsional capacity, i.e. $\underline h_{\mathcal M_0}((0,b);y_0)=0$.
\end{remark}

\begin{theorem}[Coarse upper bound]
\label{prop:coarse_upper_atomic}
    The supremum $\overline h_{\mathcal M_0}(q;y_0)$ is attained by a measure $\nu_N^+\in\Pi_{\mathcal M_0}(y_0)$ supported on at most three points, i.e.,
    \begin{align}
        \nu_N^+
        =
        \sum_{j=1}^{s}\lambda_j\delta_{x_j},
        \qquad s\leq 3,
    \end{align}
    where $x_j\in P$, $\lambda_j\geq 0$, and
    \begin{align}
        \sum_{j=1}^{s}\lambda_j=f_N,
        \qquad
        \sum_{j=1}^{s}\lambda_jx_j=f_Nc.
    \end{align}
\end{theorem}

\begin{proof}
    Since $P$ is compact and $f_N$ is fixed, $\Pi_{\mathcal M_0}(y_0)$ is weakly compact.
    It is also nonempty, since it contains $\nu_N^\star$. For fixed $q$ and $c$, the function
    $x\mapsto\|a+bJ(x-c)\|_2$ is continuous on $P$.
    Therefore, the objective is a continuous linear
    functional of $\nu_N$, and its supremum is attained
    at an extreme point of $\Pi_{\mathcal M_0}(y_0)$.    The feasible set is defined by three scalar
    moment constraints: one zeroth-moment and two
    first-moment constraints. By the extreme-point
    characterization of moment sets
    \cite{winkler1988extreme}, its extreme points
    are supported on at most three points.
\end{proof}

\begin{corollary}[Polygonal contact patch]
\label{cor:coarse_upper_polygon}
    Under the assumptions of
    Theorem~\ref{prop:coarse_upper_atomic},
    suppose additionally that $P$ is a convex polygon
    with vertices $v_1,\ldots,v_K$.

    Then the coarse upper bound is given exactly by
    the linear program
    \begin{align}
    \boxed{
    \begin{aligned}
        \overline h_{\mathcal M_0}(q;y_0)
        =
        \mu\max_{\lambda_1,\ldots,\lambda_K}\quad&
        \sum_{k=1}^K
        \lambda_k\|a+bJ(v_k-c)\|_2\\
        \text{s.t.}\quad&
        \lambda_k\geq0,\qquad k=1,\ldots,K,\\
        &
        \sum_{k=1}^K\lambda_k=f_N,\\
        &
        \sum_{k=1}^K\lambda_kv_k=f_Nc.
    \end{aligned}}
    \end{align}
    Moreover, an optimal solution exists with
    at most three nonzero weights $\lambda_k$.
\end{corollary}

\begin{proof}
    By Theorem~\ref{prop:coarse_upper_atomic},
    an optimal measure can be chosen as
    \begin{align}
        \nu_N^+=\sum_{j=1}^{s}\lambda_j\delta_{x_j},
        \qquad s\leq3.
    \end{align}
    Since $P$ is a convex polygon, each $x_j\in P$ admits a convex decomposition
    \begin{align}
        x_j=\sum_{k=1}^K\alpha_{jk}v_k,
        \qquad
        \alpha_{jk}\geq0,
        \qquad
        \sum_{k=1}^K\alpha_{jk}=1.
    \end{align}
    By convexity of $x\mapsto\|a+bJ(x-c)\|_2$,
    \begin{align}
        \|a+bJ(x_j-c)\|_2
        \leq
        \sum_{k=1}^K
        \alpha_{jk}\|a+bJ(v_k-c)\|_2.
    \end{align}
    Hence, replacing each mass $\lambda_j$ at $x_j$ by masses $\lambda_j\alpha_{jk}$ at the vertices
    preserves the zeroth and first moments without decreasing the objective.
    Therefore, an optimal measure supported on the vertices exists, and the upper bound is given
    by the stated LP. Finally, the LP has three equality constraints. Thus, an optimal basic feasible solution exists with at most three nonzero weights.
\end{proof}


\section{$\mathcal{M}_F$ case (taxel forces without CoP)}
\label{sec:MF}

We now assume that the contact patch is partitioned into $r$ pairwise disjoint
Borel measurable sets $P_1,\ldots,P_r$, or \textit{taxels}, such that
$P=\bigcup_{i=1}^r P_i$.

\begin{definition}[Taxel force measurement]\label{def:MF}
    Given the true unknown normal force measure $\nu_N^\star$ and the partition $\{P_i\}_{i=1}^r$, the \textit{taxel force measurement} is defined as
    \begin{align}
    y_F = \mathcal M_F(\nu_N^\star) \triangleq
    \left(
    \nu_N^\star(P_1),\ldots,\nu_N^\star(P_r)
    \right).
    \end{align}
    Equivalently, its measurement functions are
    $\phi_i(x)=\mathbf 1_{P_i}(x)$, $i=1,\ldots,r$.
    The measured taxel forces are denoted  $F_i\triangleq\nu_N^\star(P_i)$, $i=1,\ldots,r$.
\end{definition}

\begin{remark}
     Since the taxels form a partition of $P$, $y_F$ determines the resultant normal force $f_N=m_0(\nu_N^\star) = \sum_{i=1}^r F_i$. However, the CoP $c[\nu_N^\star]$ is not determined by $y_F$.
\end{remark}

\begin{proposition}[Compatible CoP set from taxel forces]
\label{prop:cop_taxel}
    Assume that each taxel $P_i$ is convex. The set CoPs compatible with a given a taxel force observation $y_F = (F_1,...,F_r)$ takes the form
    \begin{align}
        \mathcal C_{\mathcal M_F}(y_F)
        =
        \left\{
        \frac{1}{f_N}
        \sum_{i=1}^r F_i c_i
        :
        c_i\in P_i,\quad i=1,\ldots,r
        \right\}.
    \end{align}
    % Equivalently, $\mathcal C_{\mathcal M_F}(y_F)$ is the weighted Minkowski sum $\mathcal C_{\mathcal{M}_F}(y_F) = \sum_{i=1}^r\frac{F_i}{f_N}P_i$. 
    and $\operatorname{diam} \mathcal C_{\mathcal M_F}(y_F)\le\frac{1}{f_N}\sum_{i=1}^r F_i\operatorname{diam}(P_i)\le\max_i\operatorname{diam}(P_i).$
\end{proposition}
\begin{proof}
  Let $c\in C_{\mathcal M_F}(y_F)$. By of $C_{\mathcal M_F}(y_F)$, there exist $\nu_N\in\Pi_{\mathcal M_F}(y_F)$ such that $c[\nu_N]=c$. By definition of of $\Pi_{\mathcal M_F}(y_F)$, we must have $\nu_N(P_i)=F_i$ for every $i=1,\ldots,r$. For each $i$ let us define
\begin{align}
    c_i \triangleq \begin{cases}
                        \frac{1}{F_i}\int_{P_i}x\,d\nu_N(x) \qquad &\text{if}~F_i >0\\
                        \text{any point in $P_i$}~ \qquad &\text{if}~ F_i=0
                    \end{cases}
\end{align}
Since $\nu_N$ is nonnegative and $P_i$ is convex, every $c_i$ as defined above must belong to $P_i$. Since the taxels form a disjoint partition of $P$, we have
\begin{align}
    c[\nu_N]
    =
    \frac{1}{f_N}\int_P x\,d\nu_N(x)
    =
    \frac{1}{f_N}\sum_{i=1}^r
    \int_{P_i}x\,d\nu_N(x)
    =
    \frac{1}{f_N}\sum_{i=1}^r F_i c_i.
\end{align}
Therefore $\mathcal C_{\mathcal M_F}(y_F)    \subseteq    \left\{    \frac{1}{f_N}\sum_{i=1}^rF_i c_i: c_i\in P_i    \right\}$.  Conversely, consider any choice of $c_i\in P_i$ and define the measure $\nu_N=\sum_{i=1}^r F_i\delta_{c_i}$. Since the taxels are pairwise disjoint, $\nu_N(P_i)=F_i$ for every $i$, and hence
$\nu_N\in\Pi_{\mathcal M_F}(y_F)$.
Moreover $
    c[\nu_N]
    =
    \frac{1}{f_N}\sum_{i=1}^rF_i c_i$. 
Thus the other inclusion, and equality. Finally, consider any two points
$c,c'\in\mathcal C_{\mathcal M_F}(y_F)$.
By the previous characterization, there exist $c_i,c_i'\in P_i$ such that
\begin{align}
    c=\frac{1}{f_N}\sum_{i=1}^rF_i c_i,
    \qquad
    c'=\frac{1}{f_N}\sum_{i=1}^rF_i c_i'.
\end{align}
By the triangle inequality,
\begin{align}
    \|c-c'\|_2
    =
    \left\|
    \frac{1}{f_N}
    \sum_{i=1}^rF_i(c_i-c_i')
    \right\|_2 \le
    \frac{1}{f_N}
    \sum_{i=1}^rF_i\|c_i-c_i'\|_2\\
    \le
    \frac{1}{f_N}
    \sum_{i=1}^rF_i\operatorname{diam}(P_i).
\end{align}
Taking the supremum over all $c,c'$ gives $\operatorname{diam}\mathcal C_{\mathcal M_F}(y_F)\le\frac{1}{f_N}\sum_{i=1}^rF_i\operatorname{diam}(P_i)$. Since $F_i\ge0$ and $\sum_iF_i=f_N$, we have  $\frac{1}{f_N}\sum_{i=1}^rF_i\operatorname{diam}(P_i)\le\max_i\operatorname{diam}(P_i)$ which completes the proof.
\end{proof}

\begin{theorem}[Exact lower bound for taxel force measurements]
\label{thm:taxel_lower}
    Let $y_F=(F_1,\ldots,F_r)=\mathcal M_F(\nu_N^\star)$, where $F_i=\nu_N^\star(P_i)\geq 0$, and assume that
    $f_N=\sum_{i=1}^r F_i>0$. Then, for every $q=(a,b)\in\mathbb R^2\times\mathbb R$,
    \begin{align}
        \label{eq:taxel_lower_reduction}
        \boxed{\underline h_{\mathcal M_F}(q;y_F)
        =
        \mu
        \inf_{\substack{c_i\in P_i\\i=1,\ldots,r}}
        \sum_{i=1}^r
        F_i\left\|a+bJ(c_i-\bar c)\right\|_2}
    \end{align}
where     \begin{align}
        \bar c
        \triangleq
        \frac{1}{f_N}\sum_{i=1}^r F_i c_i.
    \end{align}
    This finite-dimensional optimization problem is convex. 
    Moreover, if each $\overline{P_i}$ is a polygon, the infimum in \eqref{eq:taxel_lower_reduction} equals the
    minimum obtained by replacing $P_i$ with $\overline{P_i}$, and the resulting problem becomes an SOCP.
\end{theorem}

\begin{proof}
    Let $\nu_N\in\Pi_{\mathcal M_F}(y_F)$. For each $i$ such that $F_i>0$, define
    \begin{align}
        c_i=\frac{1}{F_i}\int_{P_i}x\,d\nu_N(x).
    \end{align}
    By convexity of $P_i$, we have $c_i\in P_i$. For $F_i=0$, choose any $c_i\in P_i$.
    Since the taxels form a partition of $P$,
    \begin{align}
        c[\nu_N]
        =
        \frac{1}{f_N}\sum_{i=1}^rF_ic_i
        =\bar c.
    \end{align}
    Applying Jensen's inequality on each taxel gives
    \begin{align}
        \int_{P_i}
        \|a+bJ(x-\bar c)\|_2\,d\nu_N(x)
        \geq
        F_i\|a+bJ(c_i-\bar c)\|_2.
    \end{align}
    Summing over all taxels yields the inequality
    \begin{align}
        \underline h_{\mathcal M_F}(q;y_F)
        \geq
        \mu\inf_{c_i\in P_i}
        \sum_{i=1}^rF_i\|a+bJ(c_i-\bar c)\|_2.
    \end{align}

    Conversely, for any choice $c_i\in P_i$, the measure
    \begin{align}
        \nu_N=\sum_{i=1}^rF_i\delta_{c_i}
    \end{align}
    belongs to $\Pi_{\mathcal M_F}(y_F)$ and satisfies
    $c[\nu_N]=\bar c$. Its support function equals
    \begin{align}
        h_{\mathcal W(\nu_N)}(q)
        =
        \mu\sum_{i=1}^r
        F_i\|a+bJ(c_i-\bar c)\|_2.
    \end{align}
    This establishes the reverse inequality and hence equality.

    Finally, $\bar c$ is affine in $(c_1,\ldots,c_r)$,
    so the objective is convex.
    If the taxel closures are polygons, the objective is
    continuous and the product of the taxels is dense in
    the product of their compact closures.
    Consequently, the infimum equals the minimum over
    $c_i\in\overline{P_i}$, which admits an SOCP formulation.
\end{proof}

\begin{theorem}[Exact upper bound for polygonal taxels]
\label{thm:taxel_upper}
    Under the assumptions of Theorem~\ref{thm:taxel_lower},
    suppose additionally that each taxel closure is a polygon $\overline{P_i} = \operatorname{conv}\{v_{i1},\ldots,v_{iK_i}\}$ where $v_{ik}\in\mathbb R^2$ are its vertices. 
    Let $\lambda_{ik} \geq 0$ such that $\sum_{k=1}^{K_i}\lambda_{ik}=F_i$, define
    \begin{align}
        \bar c(\lambda)
        \triangleq
        \frac{1}{f_N}
        \sum_{i=1}^r\sum_{k=1}^{K_i}
        \lambda_{ik}v_{ik}.
    \end{align}
    Then, for every $q=(a,b)\in\mathbb R^2\times\mathbb R$,
    \begin{align}
        \label{eq:taxel_upper_reduction}
        \boxed{\overline h_{\mathcal M_F}(q;y_F)
        =
        \mu
        \max_{\substack{
            \lambda_{ik}\geq0\\
            \sum_{k=1}^{K_i}\lambda_{ik}=F_i,\ \forall i
        }}
        \sum_{i=1}^r\sum_{k=1}^{K_i}
        \lambda_{ik}
        \left\|
        a+bJ\bigl(v_{ik}-\bar c(\lambda)\bigr)
        \right\|_2}
    \end{align}
    This is a finite-dimensional optimization problem whose objective is generally nonconcave.
    If some vertices $v_{ik}$ do not belong to their corresponding taxels $P_i$, the equality holds for the supremum value, although the maximizing vertex-supported measure need not be admissible.
\end{theorem}

\begin{proof}
    Let $\nu_N\in\Pi_{\mathcal M_F}(y_F)$ and
    $c=c[\nu_N]$.

    Since the function
    $x\mapsto\|a+bJ(x-c)\|_2$ is convex,
    any normal force distribution on $P_i$ can be
    replaced by a nonnegative distribution on the
    vertices of $\overline{P_i}$ that preserves its
    total force and first moment, without decreasing
    the integral of this function.

    Applying this construction to every taxel
    produces weights $\lambda_{ik}\geq0$ satisfying
    \begin{align}
        \sum_{k=1}^{K_i}\lambda_{ik}=F_i,
        \qquad
        \bar c(\lambda)=c,
    \end{align}
    such that
    \begin{align}
        h_{\mathcal W(\nu_N)}(q)
        \leq
        \mu\sum_{i=1}^r\sum_{k=1}^{K_i}
        \lambda_{ik}
        \|a+bJ(v_{ik}-\bar c(\lambda))\|_2.
    \end{align}
    Hence the original supremum is bounded above by
    \eqref{eq:taxel_upper_reduction}.

    Conversely, consider any feasible weights
    $\lambda_{ik}$. Since
    $v_{ik}\in\overline{P_i}$, each vertex can be
    approximated by a sequence of points in $P_i$.
    Assigning the force $\lambda_{ik}$ to these
    points gives a sequence of measures compatible
    with $y_F$ whose CoPs converge to $\bar c(\lambda)$
    and whose support values converge to the
    objective in \eqref{eq:taxel_upper_reduction}.

    Therefore, the original supremum is at least the
    maximum in \eqref{eq:taxel_upper_reduction},
    establishing equality.
\end{proof}

Contrary to the coarse measurement $y_0$, the taxel force measurement $y_F$ does not determine the CoP. The bounds therefore retain the bi-level structure


Thus, although $\mathcal M_F$ does not directly measure the CoP, the taxel
resultants constrain it. In particular, finer spatial measurements reduce the
set of compatible CoPs.

\begin{proposition}[Tightening under taxel refinement without CoP]
\label{prop:taxel_refinement_MF}
Let $\mathcal M_F'$ denote a taxel-force measurement associated with a
partition that refines the partition associated with $\mathcal M_F$, and let
\begin{align}
y_F=\mathcal M_F(\nu_N^\star),
\qquad
y_F'=\mathcal M_F'(\nu_N^\star).
\end{align}
Then
\begin{align}
\Pi_{\mathcal M_F'}(y_F')
&\subseteq
\Pi_{\mathcal M_F}(y_F),\\
\mathcal C_{\mathcal M_F'}(y_F')
&\subseteq
\mathcal C_{\mathcal M_F}(y_F).
\end{align}
Consequently, for every $q\in\mathbb R^2\times\mathbb R$,
\begin{align}
\boxed{
\underline h_{\mathcal M_F}(q;y_F)
\le
\underline h_{\mathcal M_F'}(q;y_F')
\le
h_{\mathcal W(\nu_N^\star)}(q)
\le
\overline h_{\mathcal M_F'}(q;y_F')
\le
\overline h_{\mathcal M_F}(q;y_F).
}
\end{align}

Moreover, for $q=(a,b)$,
\begin{align}
\boxed{
0
\le
\overline h_{\mathcal M_F}(q;y_F)
-
\underline h_{\mathcal M_F}(q;y_F)
\le
2\mu |b|
\sum_{i=1}^r
F_i\operatorname{diam}(P_i).
}
\end{align}
In particular,
\begin{align}
\overline h_{\mathcal M_F}(q;y_F)
-
\underline h_{\mathcal M_F}(q;y_F)
\le
2\mu |b|f_N
\max_i\operatorname{diam}(P_i).
\end{align}
Hence, along a sequence of refining taxel partitions for which
$\max_i\operatorname{diam}(P_i)\rightarrow0$, the lower and upper
measurement-conditioned bounds converge to
$h_{\mathcal W(\nu_N^\star)}(q)$.
\end{proposition}

The monotonicity above is associated with \textit{partition refinement}, not
with the number of taxels $r$ alone.

\paragraph{Fixed-reference approximation.}
A simpler approximation is obtained by replacing the unknown CoP by a fixed
estimate $\hat c$ and using $\hat c$ only as the wrench reference, without
enforcing the CoP consistency constraint
$\int_P(x-\hat c)\,d\nu_N=0$. In that case, the taxel constraints are
independent and the optimization decouples:
\begin{align}
\underline h_F^{\mathrm{ref}}(q;y_F,\hat c)
&=
\mu\sum_{i=1}^r
F_i
\inf_{x\in P_i}
\|a+bJ(x-\hat c)\|_2,\\
\overline h_F^{\mathrm{ref}}(q;y_F,\hat c)
&=
\mu\sum_{i=1}^r
F_i
\sup_{x\in P_i}
\|a+bJ(x-\hat c)\|_2.
\end{align}
These expressions are exact for this fixed-reference decoupled approximation,
but not for the original $\mathcal M_F$-conditioned problem. In particular,
fixing $\hat c$ while retaining
$\int_P(x-\hat c)\,d\nu_N=0$ would still couple the taxels.


\section{$\mathcal{M}_T$ case (taxel forces with CoP)}
\label{sec:MT}

\begin{definition}[Combined coarse and taxel force measurement]
If both coarse and taxel force information are available, we define the \textit{combined coarse-and-taxel measurement}
\begin{align}
    y_T=\mathcal M_T(\nu_N^\star)
    &\triangleq
    \bigl(
    \mathcal M_0(\nu_N^\star),
    \mathcal M_F(\nu_N^\star)
    \bigr) \\
    &= \left(
    (f_N,f_Nc[\nu_N^\star]), F_1,\ldots,F_r\right).
\end{align}
\end{definition}
Since $y_T$ determines the CoP $c^\star=c[\nu_N^\star]$, the set of compatible CoPs is reduced to a singleton $\mathcal C_{\mathcal M_T}(y_T)=\{c^\star\}$. Hence, the outer optimization therefore disappears from \eqref{eq:parametric_gmp_lower},\eqref{eq:parametric_gmp_upper} and the $\mathcal M_T$-conditioned bounds problem reduce to the GMPs 
    \begin{align}
    \underline h_{\mathcal M_T}(q;y_T)
    =
    \inf_{\substack{
    \nu_N(P_i)=F_i,\ i=1,\ldots,r\\
    \int_P(x-c^\star)\,d\nu_N=0}}
    \mu\int_P
    \|a+bJ(x-c^\star)\|_2\,d\nu_N(x),
\end{align}
\begin{align}
    \overline h_{\mathcal M_T}(q;y_T)
    =
    \sup_{\substack{
    \nu_N(P_i)=F_i,\ i=1,\ldots,r\\
    \int_P(x-c^\star)\,d\nu_N=0}}
    \mu\int_P
    \|a+bJ(x-c^\star)\|_2\,d\nu_N(x).
\end{align}
The taxel force constraints are disjoint, but the CoP consistency constraint $\int_P(x-c^\star)\,d\nu_N=0$ couples the measure $\nu_N$ across the taxels.

\begin{theorem}[Exact lower bound for combined measurements]
\label{thm:combined_lower}
    Let $P\subset\mathbb R^2$ be compact and let
    $(P_i)_{i=1}^r$ be a partition of $P$ into nonempty,
    pairwise disjoint, convex Borel sets.

    Given the combined observation
    \begin{align}
        y_T
        &=
        \mathcal M_T(\nu_N^\star)
        =(y_0,y_F),\\
        y_0&=(f_N,f_Nc),
        \qquad
        y_F=(F_1,\ldots,F_r),
    \end{align}
    where $f_N=\sum_{i=1}^rF_i>0$, the measurement-conditioned
    lower bound satisfies
    \begin{align}
        \label{eq:combined_lower_reduction}
        \underline h_{\mathcal M_T}(q;y_T)
        =
        \mu
        \inf_{\substack{
            c_i\in P_i,\ i=1,\ldots,r\\
            \sum_{i=1}^rF_ic_i=f_Nc
        }}
        \sum_{i=1}^r
        F_i\left\|a+bJ(c_i-c)\right\|_2.
    \end{align}
    The resulting finite-dimensional optimization problem is convex.

    Moreover, if each $\overline{P_i}$ is a polygon and
    $c\in\operatorname{ri}\mathcal C_{\mathcal M_F}(y_F)$,
    the infimum equals the minimum obtained by replacing
    $P_i$ with $\overline{P_i}$, and the resulting problem
    admits an SOCP formulation.
\end{theorem}


\begin{theorem}[Exact upper bound for combined measurements]
\label{thm:combined_upper}
    Under the assumptions of Theorem~\ref{thm:combined_lower},
    suppose additionally that each taxel closure is a polygon,
    \begin{align}
        \overline{P_i}
        =
        \operatorname{conv}
        \{v_{i1},\ldots,v_{iK_i}\},
    \end{align}
    where $v_{ik}\in\mathbb R^2$ are its vertices.

    If $c\in\operatorname{ri}\mathcal C_{\mathcal M_F}(y_F)$,
    then, for every $q=(a,b)\in\mathbb R^2\times\mathbb R$,
    \begin{align}
        \label{eq:combined_upper_reduction}
        \overline h_{\mathcal M_T}(q;y_T)
        =
        \mu\max_{\lambda_{ik}\geq0}\quad&
        \sum_{i=1}^r\sum_{k=1}^{K_i}
        \lambda_{ik}
        \left\|a+bJ(v_{ik}-c)\right\|_2\\
        \text{subject to}\quad&
        \sum_{k=1}^{K_i}\lambda_{ik}=F_i,
        \qquad i=1,\ldots,r,\\
        &
        \sum_{i=1}^r\sum_{k=1}^{K_i}
        \lambda_{ik}v_{ik}=f_Nc.
    \end{align}
    The resulting finite-dimensional optimization problem
    is a linear program.

    The equality concerns the supremum value; if some
    vertices do not belong to their corresponding taxels,
    the maximizing vertex-supported measure need not be
    admissible.
\end{theorem}

\begin{proposition}[Refinement from coarse and taxel-force measurements]
\label{prop:taxel_refinement}
The compatible measure family associated with the combined measurement is
\begin{align}
\Pi_{\mathcal M_T}(y_T)
=
\Pi_{\mathcal M_0}(y_0)
\cap
\Pi_{\mathcal M_F}(y_F).
\end{align}
Consequently, for every wrench direction
$q\in\mathbb R^2\times\mathbb R$,
\begin{align}
\boxed{
\max\left\{
\underline h_{\mathcal M_0}(q;y_0),
\underline h_{\mathcal M_F}(q;y_F)
\right\}
\le
\underline h_{\mathcal M_T}(q;y_T)
\le
h_{\mathcal W(\nu_N^\star)}(q)
}
\end{align}
and
\begin{align}
\boxed{
h_{\mathcal W(\nu_N^\star)}(q)
\le
\overline h_{\mathcal M_T}(q;y_T)
\le
\min\left\{
\overline h_{\mathcal M_0}(q;y_0),
\overline h_{\mathcal M_F}(q;y_F)
\right\}.
}
\end{align}
Thus, $\mathcal M_T$ provides bounds that are at least as tight as those
obtained from either $\mathcal M_0$ or $\mathcal M_F$ alone, while
$\mathcal M_0$ and $\mathcal M_F$ are not in general ordered.
\end{proposition}

The exact $\mathcal M_T$-conditioned bounds retain both the taxel-force
constraints and the global first-moment constraint associated with the measured
CoP. Simpler bounds can be obtained by treating the taxels independently while
using the same measured $c^\star$ and $F_i$.

\begin{proposition}[Constructive cellwise support bounds]
Let $q=(a,b)\in\mathbb R^2\times\mathbb R$, and define, for
$i=1,\ldots,r$,
\begin{align}
g_i^{\min}(q;c^\star)
&\triangleq
\inf_{x\in P_i}
\|a+bJ(x-c^\star)\|_2,\\
g_i^{\max}(q;c^\star)
&\triangleq
\sup_{x\in P_i}
\|a+bJ(x-c^\star)\|_2.
\end{align}
Then, for every $\nu_N\in\Pi_{\mathcal M_T}(y_T)$,
\begin{align}
\boxed{
\mu\sum_{i=1}^r
F_i g_i^{\min}(q;c^\star)
\le
h_{\mathcal W(\nu_N)}(q)
\le
\mu\sum_{i=1}^r
F_i g_i^{\max}(q;c^\star).
}
\end{align}
\end{proposition}

\begin{proof}
For every $x\in P_i$,
\begin{align}
g_i^{\min}(q;c^\star)
\le
\|a+bJ(x-c^\star)\|_2
\le
g_i^{\max}(q;c^\star).
\end{align}
Since $\nu_N$ is nonnegative and
$\nu_N(P_i)=F_i$ for every
$\nu_N\in\Pi_{\mathcal M_T}(y_T)$, integration over $P_i$ gives
\begin{align}
F_i g_i^{\min}(q;c^\star)
\le
\int_{P_i}
\|a+bJ(x-c^\star)\|_2\,d\nu_N(x)
\le
F_i g_i^{\max}(q;c^\star).
\end{align}
Summing over the pairwise disjoint taxels and using
$P=\bigcup_{i=1}^rP_i$ yields
\begin{align}
\sum_{i=1}^rF_i g_i^{\min}(q;c^\star)
\le
\int_P
\|a+bJ(x-c^\star)\|_2\,d\nu_N(x)
\le
\sum_{i=1}^rF_i g_i^{\max}(q;c^\star).
\end{align}
The result follows from Proposition~\ref{prop:support}.
\end{proof}

Consequently,
\begin{align}
\boxed{
\mu\sum_{i=1}^r F_i g_i^{\min}(q;c^\star)
\le
\underline h_{\mathcal M_T}(q;y_T)
\le
\overline h_{\mathcal M_T}(q;y_T)
\le
\mu\sum_{i=1}^r F_i g_i^{\max}(q;c^\star).
}
\end{align}
These cellwise bounds are directly computable from the taxel forces, the
measured CoP, and the taxel geometry, but need not be sharp. In particular,
independently attaining the minimum or maximum within each taxel may violate
the global first-moment constraint contained in $\mathcal M_T$.

\begin{proposition}[Resolution-dependent tightening with CoP]
\label{prop:taxel_resolution_MT}
For every $q=(a,b)\in\mathbb R^2\times\mathbb R$,
\begin{align}
\boxed{
0
\le
\overline h_{\mathcal M_T}(q;y_T)
-
\underline h_{\mathcal M_T}(q;y_T)
\le
\mu |b|
\sum_{i=1}^r
F_i\operatorname{diam}(P_i).
}
\end{align}
In particular,
\begin{align}
\overline h_{\mathcal M_T}(q;y_T)
-
\underline h_{\mathcal M_T}(q;y_T)
\le
\mu |b| f_N
\max_i\operatorname{diam}(P_i).
\end{align}

Moreover, if a taxel partition is refined while retaining the same measured
CoP $c^\star$, the compatible measure family can only shrink and therefore the
lower bound can only increase while the upper bound can only decrease. Along
any sequence of refining partitions satisfying
\begin{align}
\max_i\operatorname{diam}(P_i)\rightarrow0,
\end{align}
we have
\begin{align}
\underline h_{\mathcal M_T}(q;y_T)
\rightarrow
h_{\mathcal W(\nu_N^\star)}(q),
\qquad
\overline h_{\mathcal M_T}(q;y_T)
\rightarrow
h_{\mathcal W(\nu_N^\star)}(q).
\end{align}
\end{proposition}

\subsection{Current issues and open questions}

The two taxel cases differ in the role played by the CoP. With
$\mathcal M_F$, the CoP remains an outer optimization variable and the
compatible CoP set $\mathcal C_{\mathcal M_F}(y_F)$ contracts as the taxel
partition is refined. With $\mathcal M_T$, the CoP is fixed by the measurement,
but its first-moment constraint couples the taxels in the inner GMP. The main
remaining questions concern efficient finite-dimensional solution of these
problems, the sharpness of the resolution-dependent bounds above, and the
additional tightening provided by the CoP constraint in $\mathcal M_T$ relative
to $\mathcal M_F$.

% % \paragraph{Open research directions.}
% % The measurement-conditioned bounds tighten monotonically as we add richer
% % measurements. It follows directly from the inclusion of the corresponding compatible
% % measure families. Interesting questions are about quantifying / identifying:
% % \begin{itemize}
% %     \item How does the uncertainty
% %     $\overline h_{\mathcal M}(q;y)-\underline h_{\mathcal M}(q;y)$
% %     decrease as the spatial resolution of the tactile measurement increases?
% %     \item Can one prove convergence of the measurement-conditioned bounds to the
% %     true support function as the taxel partition is refined, and obtain a rate in
% %     terms of the cell diameter?
% %     \item For a given wrench direction $q$, which measurements of the normal force
% %     measure are sufficient to identify $h_{\mathcal W(\nu_N)}(q)$ exactly?
% %     \item Which wrench directions remain fundamentally unidentifiable from coarse
% %     measurements such as resultant normal force and center of pressure?
% %     \item What is the minimal spatial information required to certify a prescribed
% %     force or torsional capacity?
% % \end{itemize}
% % \subsection{Pure torsion as a special case}

\begin{thebibliography}{9}

\bibitem{evans2015measure}
L.~C. Evans and R.~F. Gariepy,
\newblock \emph{Measure Theory and Fine Properties of Functions},
\newblock Revised ed., CRC Press, Boca Raton, FL, 2015.

% \bibitem{lasserre2010moments}
% J.-B. Lasserre,
% \newblock \emph{Moments, Positive Polynomials and Their Applications},
% \newblock Imperial College Press Optimization Series, vol.~1,
% Imperial College Press, London, 2010.

% \bibitem{winkler1988extreme}
% G.~Winkler,
% \newblock ``Extreme points of moment sets,''
% \newblock \emph{Mathematics of Operations Research},
% vol.~13, no.~4, pp.~581--587, 1988.

\end{thebibliography}
\newpage
\section*{APPENDIX}
\begin{proposition}[$\mathcal M_0(\nu_N)$ does not determine $\mathcal W(\nu_N)$]
    There exist a compact convex contact patch $P\subset\mathbb R^2$ and two
    distinct normal force measures $\nu_N^{(1)},\nu_N^{(2)}\in\mathcal M_+(P)$ such that
    \begin{align}
    \mathcal M_0(\nu_N^{(1)}) = \mathcal M_0(\nu_N^{(2)}) ~~\text{and}~~\mathcal W(\nu_N^{(1)})\neq\mathcal W(\nu_N^{(2)}).
    \end{align}
\end{proposition}
    \begin{proof}
    For any $L>0$ and $f_N>0$, consider the segment $P=[0,L]\times\{0\}\subset\mathbb R^2$ with midpoint $c=\left(\frac{L}{2},0\right)$ and define the two normal force measures
    \begin{align}
    \nu_N^{(1)}
    &=
    f_N\delta_c,\\
    \nu_N^{(2)}
    &=
    \frac{f_N}{2}\delta_{(0,0)}
    +
    \frac{f_N}{2}\delta_{(L,0)}.
    \end{align}
    The two measures the same resultant normal force and CoP :
    \begin{align}
    \mathcal M_0(\nu_N^{(1)}) = \mathcal M_0(\nu_N^{(2)}) = \left( f_N,f_N\left(\frac{L}{2},0\right)\right).
    \end{align}
    But using Proposition~\ref{prop:support} in the pure torsion support direction $q=(0,1)$, we have 
    \begin{align}
    h_{\mathcal W(\nu_N^{(1)})}(0,1)
    =
    0 ~~~\neq~~~
    h_{\mathcal W(\nu_N^{(2)})}(0,1)
    =
    \mu f_N\frac{L}{2}.
    \end{align}
    Hence $\mathcal W(\nu_N^{(1)}) \neq \mathcal W(\nu_N^{(2)})$. \textcolor{red}{TO CHECK with Yann}
\end{proof}


\begin{proposition}[Compatible CoP set from taxel forces]
\label{prop:cop_taxel}
    Assume that each taxel $P_i$ is convex. The set CoPs compatible with a given a taxel force observation $y_F = (F_1,...,F_r)$ takes the form
    \begin{align}
        \mathcal C_{\mathcal M_F}(y_F)
        =
        \left\{
        \frac{1}{f_N}
        \sum_{i=1}^r F_i c_i
        :
        c_i\in P_i,\quad i=1,\ldots,r
        \right\}.
    \end{align}
    % Equivalently, $\mathcal C_{\mathcal M_F}(y_F)$ is the weighted Minkowski sum $\mathcal C_{\mathcal{M}_F}(y_F) = \sum_{i=1}^r\frac{F_i}{f_N}P_i$. 
    and $\operatorname{diam} \mathcal C_{\mathcal M_F}(y_F)\le\frac{1}{f_N}\sum_{i=1}^r F_i\operatorname{diam}(P_i)\le\max_i\operatorname{diam}(P_i).$
\end{proposition}
\begin{proof}
  Let $c\in C_{\mathcal M_F}(y_F)$. By of $C_{\mathcal M_F}(y_F)$, there exist $\nu_N\in\Pi_{\mathcal M_F}(y_F)$ such that $c[\nu_N]=c$. By definition of of $\Pi_{\mathcal M_F}(y_F)$, we must have $\nu_N(P_i)=F_i$ for every $i=1,\ldots,r$. For each $i$ let us define
\begin{align}
    c_i \triangleq \begin{cases}
                        \frac{1}{F_i}\int_{P_i}x\,d\nu_N(x) \qquad &\text{if}~F_i >0\\
                        \text{any point in $P_i$}~ \qquad &\text{if}~ F_i=0
                    \end{cases}
\end{align}
Since $\nu_N$ is nonnegative and $P_i$ is convex, every $c_i$ as defined above must belong to $P_i$. Since the taxels form a disjoint partition of $P$, we have
\begin{align}
    c[\nu_N]
    =
    \frac{1}{f_N}\int_P x\,d\nu_N(x)
    =
    \frac{1}{f_N}\sum_{i=1}^r
    \int_{P_i}x\,d\nu_N(x)
    =
    \frac{1}{f_N}\sum_{i=1}^r F_i c_i.
\end{align}
Therefore $\mathcal C_{\mathcal M_F}(y_F)    \subseteq    \left\{    \frac{1}{f_N}\sum_{i=1}^rF_i c_i: c_i\in P_i    \right\}$.  Conversely, consider any choice of $c_i\in P_i$ and define the measure $\nu_N=\sum_{i=1}^r F_i\delta_{c_i}$. Since the taxels are pairwise disjoint, $\nu_N(P_i)=F_i$ for every $i$, and hence
$\nu_N\in\Pi_{\mathcal M_F}(y_F)$.
Moreover $
    c[\nu_N]
    =
    \frac{1}{f_N}\sum_{i=1}^rF_i c_i$. 
Thus the other inclusion, and equality. Finally, consider any two points
$c,c'\in\mathcal C_{\mathcal M_F}(y_F)$.
By the previous characterization, there exist $c_i,c_i'\in P_i$ such that
\begin{align}
    c=\frac{1}{f_N}\sum_{i=1}^rF_i c_i,
    \qquad
    c'=\frac{1}{f_N}\sum_{i=1}^rF_i c_i'.
\end{align}
By the triangle inequality,
\begin{align}
    \|c-c'\|_2
    =
    \left\|
    \frac{1}{f_N}
    \sum_{i=1}^rF_i(c_i-c_i')
    \right\|_2 \le
    \frac{1}{f_N}
    \sum_{i=1}^rF_i\|c_i-c_i'\|_2\\
    \le
    \frac{1}{f_N}
    \sum_{i=1}^rF_i\operatorname{diam}(P_i).
\end{align}
Taking the supremum over all $c,c'$ gives $\operatorname{diam}\mathcal C_{\mathcal M_F}(y_F)\le\frac{1}{f_N}\sum_{i=1}^rF_i\operatorname{diam}(P_i)$. Since $F_i\ge0$ and $\sum_iF_i=f_N$, we have  $\frac{1}{f_N}\sum_{i=1}^rF_i\operatorname{diam}(P_i)\le\max_i\operatorname{diam}(P_i)$ which completes the proof.
\end{proof}
\end{document}